\begin{table}[t]
\centering\small
\setlength{\tabcolsep}{4.5pt}
\renewcommand{\arraystretch}{1.08}
\begin{tabularx}{\linewidth}{@{}Xlrrrr@{}}
\toprule
\textbf{Benchmark} & \textbf{Role} & \textbf{Direct} & \textbf{Loop} & \textbf{Gain} & \textbf{W/T/L} \\
\midrule
DTI: TAPB & Dev. replay & 93.63 & \textbf{94.02} & \textbf{+0.39} & 2/1/0 \\
Proteomics: efficacy & Held-out & \textbf{37.47} & \textbf{37.47} & +0.00 & 0/3/0 \\
Cell: VCC single-gene & Held-out & \textbf{20.01} & \textbf{20.01} & +0.00 & 0/3/0 \\
Cell: Norman double-gene & Held-out & 11.48 & \textbf{22.32} & \textbf{+10.84} & 3/0/0 \\
Cell: Tahoe drug & Held-out & 20.90 & \textbf{22.89} & \textbf{+1.99} & 1/2/0 \\
\bottomrule
\end{tabularx}
\caption{Compute-matched feedback isolation, averaged over three seeds and reported in percentage points. Held-out arms each use 800 full-client rounds and 160 development measurements. Gain is loop minus direct; W/T/L counts seed-level wins, ties and losses. The held-out total is 4/8/0, with Norman 95\% CI [4.54, 17.73] and Tahoe CI [$-$1.00, 7.46]. DTI reports a conservative development-trajectory replay; no DTI held-out claim is made.}
\label{tab:strong-effects}
\end{table}
